% Included by arbitrary_plumbing_superconformal_blocks.tex.
% Current BPZ and plumbing convention.
\section{Conventions and reproduction of the calculation}
\label{app:conventions}

This appendix supplements the block definitions above with the conventions
and numerical details needed to reproduce the calculation. We use the
PBW states and physical three-point normalizations of the long
paper, but every pairing and descendant form uses the single BPZ
convention fixed below. The algebraic double-Virasoro embedding,
oscillator realization, and reflection rules are unchanged.

The remaining choices specify the bilinear pairing, contour branches,
primary scales used by the code, and numerical evaluation. They
apply to both ordinary and inserted blocks in the $1/z$ BPZ convention.
``Norm'' always means a bilinear BPZ pairing. Generic representation
parameters are assumed unless a limit is specified.

\subsection{SCA normalization and momentum inputs}
\label{app:algebra}

The physical central charge is $c$, with
\begin{align}
 [L_m,L_n]&=(m-n)L_{m+n}
       +\frac{c}{12}(m^3-m)\delta_{m+n,0},\nonumber\\
 [L_m,G_r]&=\left(\frac m2-r\right)G_{m+r},\nonumber\\
 \{G_r,G_s\}&=2L_{r+s}
       +\frac c3\left(r^2-\frac14\right)\delta_{r+s,0}.
 \label{eq:app-sca}
\end{align}
Here $m,n\in\mathbb Z$, $r,s\in\mathbb Z+\nu$, and
$\nu=\frac12$ in NS and $\nu=0$ in R. Thus $G_0^2=L_0-c/24$.
The auxiliary algebra is the one in the paper's Sec.~3. With the
auxiliary factor first, $G$ acting on an auxiliary state of parity
$\epsilon'$ acquires $(-1)^{\epsilon'}$.
All positive modes annihilate the corresponding highest-weight states;
Ramond zero modes act on the grounds and are not separate PBW generators.

The code takes $b$ and $P$ as inputs and uses the parameter relations
in the paper's Sec.~3, fixing the remaining sign by $\beta=P/\sqrt2$.
It does not reconstruct $b$ from $c$; the embedding requires
$b\ne0$ and $b^2\ne1$. The real and imaginary momentum branches are
specified in the paper's Appendix~A. For a general complex input,
choose $P$ once and keep it at every occurrence of that module.
Reflection transports the ground-state frame as specified there;
it is not an independent square-root choice at each vertex.

\subsection{PBW order and bilinear BPZ pairing}
\label{app:pbw}

Use the physical PBW order of the paper's Sec.~2 and auxiliary words
$\Psi_{-\mathsf A}=\psi_{-s_1}\cdots\psi_{-s_t}$ with
$s_1>\cdots>s_t>0$. Repeated odd modes are reduced using
$G_{-r}^2=L_{-2r}$ and $\psi_{-s}^2=0$.
An operator product acts on the ground from right to left.
For ground labels written as signs, $|+|=0$ and $|-|=1$.

The BPZ transpose is \emph{linear and graded}:
\begin{equation}
 L_n^\dagger=L_{-n},\qquad G_r^\dagger=iG_{-r},\qquad
 \psi_r^\dagger=-i\psi_{-r},\qquad
 (XY)^\dagger=(-1)^{|X||Y|}Y^\dagger X^\dagger.
 \label{eq:app-adjoint}
\end{equation}
It does not conjugate $i$, $P$, or any other coefficient. The
bilinear pairing satisfies
\begin{equation}
 \langle\!\langle Ax|y\rangle
 =(-1)^{|A||x|}\langle\!\langle x|A^\dagger y\rangle.
 \label{eq:app-adjoint-pair}
\end{equation}
For the paper's PBW word containing $p$ supercurrents, this gives
\begin{equation}
 \langle\!\langle\mathbf L_{-A}g|y\rangle
 =i^p(-1)^{p|g|+p(p-1)/2}
 \langle\!\langle g|
 G_{r_p}\cdots G_{r_1}L_{n_\ell}\cdots L_{n_1}y\rangle.
 \label{eq:app-pbw-transpose}
\end{equation}
Adjointness of the zero modes fixes the ground metrics:
\begin{equation}
 \mathbb G_\beta\big|_0=\operatorname{diag}(1,-1),\qquad
 G_{\F,\R}\big|_0=\operatorname{diag}(1,-i).
 \label{eq:app-ground-metrics}
\end{equation}
Indeed the graded pairing gives
$e^{-i\pi/4}\langle\!\langle w^-|w^-\rangle
=i e^{i\pi/4}\langle\!\langle w^+|w^+\rangle$.
All mixed ground pairings vanish. The auxiliary oscillator metric is
diagonal; for a word with $m$ distinct modes and Ramond ground
$u^{\mathsf a}$ its entry is
$(-i)^{m+\mathsf a}(-1)^{m\mathsf a+m(m-1)/2}$.
Its inverse is the reciprocal. Set $\mathsf a=0$ in NS.

On the enlarged state space the pairing is graded:
\begin{equation}
 \langle\!\langle a\otimes x|b\otimes y\rangle
 =(-1)^{|a||x|}\langle\!\langle a|b\rangle_{\F}
  \langle\!\langle x|y\rangle_{\mathsf{SCA}}.
 \label{eq:app-product-metric}
\end{equation}
Compute a physical Gram matrix by commuting the positive modes of the
bra through the ket using \eqref{eq:app-sca}, ending on
\eqref{eq:app-ground-metrics}. Invert separately at each level and
parity. Upper Gram indices in the sewing formulas denote this inverse,
including the Ramond ground metric. Multiplying its entries by another
ground phase would count that phase twice.

Some oscillator routines use a numerical Ramond basis rather than
$w^\pm$. Any such basis change acts on both the Gram and local vertex;
it must not be attached as an extra sewing or three-point phase.
The BPZ-native direct-PBW path checks the returned coefficients
without that intermediate basis conversion.

\subsection{Local coordinates and spin frames}
\label{app:local-data}

The ordered punctures use local coordinates $(1/z,z-1,z)$.
At the first slot the BPZ inversion and its lift are
\begin{equation}
 I:z\mapsto1/z,\qquad \lambda_I(z)=-i/z.
 \label{eq:app-bpz-map}
\end{equation}
At an edge they are
\begin{equation}
 I_q:z\mapsto q/z,\qquad
 \lambda_{I_q}(z)=-i\sqrt q/z,
 \qquad \sqrt{-1}=i.
 \label{eq:app-edge-lift}
\end{equation}
The edge map is BPZ inversion followed by scaling by $q$; its
state propagator is $q^{L_0}$ and contributes no extra phase.
The fixed phases in the square-root coordinates are part of the
spin frames; $\eta_e$ records the remaining sign.
Moving punctures or permuting slots requires transporting these frames
and the fields' conformal factors, as well as their parity signs.
Exchanging weight labels alone is not a slot permutation.

Use the physical and auxiliary ground values in the paper's Secs.~2
and~4, respectively. Together with the Ward identities below, they fix
all descendant forms. In particular the factors $+i\eta$ and $i$ in
those ground values must be kept with the stated contour branch.

\subsection{Ward identities with all contour signs displayed}
\label{app:wards}

Expand the currents as $T(z)=\sum_nL_nz^{-n-2}$,
$G(z)=\sum_rG_rz^{-r-3/2}$, and
$\psi(z)=\sum_r\psi_rz^{-r-1/2}$.
For integer $m,n$ and any of the forms, the Virasoro identity is
\begin{align}
 &\sum_{j\ge0}(-1)^j\binom nj
       \rho(L_{j-m-n+1}x_1,x_2,x_3)\nonumber\\
 &\quad=\sum_{j\ge0}\binom mj\rho(x_1,L_{n+j-1}x_2,x_3)
   +(-1)^n\sum_{j\ge0}(-1)^j\binom nj
                         \rho(x_1,x_2,L_{m+j-1}x_3).
 \label{eq:app-vir-ward}
\end{align}
Here $x_i$ is a state, not a coordinate. The generalized binomial means
$\binom{x}{j}=x(x-1)\cdots(x-j+1)/j!$, with $\binom{x}{0}=1$.
Each sum terminates on finite-level input because sufficiently positive
modes annihilate a state. Negative upper binomial indices do not justify
truncating a sum prematurely.

In NS--R--R choose
\begin{equation}
 \sqrt{z(z-1)}=
 \begin{cases}
 z\sqrt{1-z^{-1}},&z\sim\infty,\\
 \sqrt{z-1}\sqrt{1+(z-1)},&z\sim1,\\
 i\sqrt z\sqrt{1-z},&z\sim0.
 \end{cases}
 \label{eq:app-contour-root}
\end{equation}
The power series roots on the right have constant coefficient one.
The last line is the upper-lip continuation between the Ramond punctures;
it is not obtained by independently taking a principal root in each
coordinate patch. Residues with multipliers
$z^m(z-1)^n\sqrt{z(z-1)}$ and
$z^m(z-1)^n/\sqrt{z(z-1)}$ give, respectively,
\begin{align}
 &\sum_{j\ge0}(-1)^j\binom{n+\frac12}{j}
 \rho_f^{(\eta)}(G_{j-m-n-\frac12}x_1,x_2,x_3)\nonumber\\
 &=i(-1)^{\epsilon_1}\sum_{j\ge0}\binom{m+\frac12}{j}
 \rho_f^{(\eta)}(x_1,G_{n+j}x_2,x_3)
 -(-1)^{\epsilon_1+f+\epsilon_2+n}
       \sum_{j\ge0}(-1)^j\binom{n+\frac12}{j}
 \rho_f^{(\eta)}(x_1,x_2,G_{m+j}x_3),
 \label{eq:app-r-g-ward}\\
 &\sum_{j\ge0}(-1)^j\binom{n-\frac12}{j}
 \rho_{\F}(\psi_{j-m-n+\frac12}x_1,x_2,x_3)\nonumber\\
 &=-i(-1)^{\epsilon'_1}\sum_{j\ge0}\binom{m-\frac12}{j}
 \rho_{\F}(x_1,\psi_{n+j}x_2,x_3)
 -(-1)^{\epsilon'_1+\epsilon'_2+n}
       \sum_{j\ge0}(-1)^j\binom{n-\frac12}{j}
 \rho_{\F}(x_1,x_2,\psi_{m+j}x_3).
 \label{eq:app-r-psi-ward}
\end{align}
The $\epsilon_1$ and $\epsilon'_1$ factors record transport past the
first state, while $\epsilon_2$ and $\epsilon'_2$ record transport past
the middle state. The additional $(-1)^f$ is the ordering sign of the
odd NS--R--R form; it is not a square-root branch change.
These formulas assume an intrinsically even NS primary at infinity.
For an intrinsically odd one, reverse the physical origin sign,
and multiply its ordered ground vertex and BPZ norm by $i$.

For all-NS vertices the multipliers contain no Ramond square root:
\begin{align}
 &\sum_{j\ge0}(-1)^j\binom nj
       \rho_a(G_{j-m-n+\frac12}x_1,x_2,x_3)\nonumber\\
 &=i(-1)^{\epsilon_1}\sum_{j\ge0}\binom mj
       \rho_a(x_1,G_{n+j-\frac12}x_2,x_3)
   +i(-1)^{\epsilon_1+\epsilon_2+n}
      \sum_{j\ge0}(-1)^j\binom nj
                      \rho_a(x_1,x_2,G_{m+j-\frac12}x_3),
 \label{eq:app-ns-g-ward}\\
 &\sum_{j\ge0}(-1)^j\binom nj
       \rho_{\F}(\psi_{j-m-n-\frac12}x_1,x_2,x_3)\nonumber\\
 &=-i(-1)^{\epsilon'_1}\sum_{j\ge0}\binom mj
       \rho_{\F}(x_1,\psi_{n+j+\frac12}x_2,x_3)
   -i(-1)^{\epsilon'_1+\epsilon'_2+n}
      \sum_{j\ge0}(-1)^j\binom nj
                      \rho_{\F}(x_1,x_2,\psi_{m+j+\frac12}x_3).
 \label{eq:app-ns-psi-ward}
\end{align}
For example, on Ramond grounds the $m=n=0$ supercurrent identity gives
\begin{align}
 \rho_f^{(\eta)}(G_{-1/2}\phi_1,w_2^\alpha,w_3^\gamma)
 &=i\rho_f^{(\eta)}(\phi_1,G_0w_2^\alpha,w_3^\gamma)\nonumber\\
 &\quad-(-1)^{f+|\alpha|}
       \rho_f^{(\eta)}(\phi_1,w_2^\alpha,G_0w_3^\gamma).
 \label{eq:app-ground-g-ward}
\end{align}
The auxiliary identity on $(\mathbf1,u^0,u^1)$ gives
$0=-i\rho_{\F}(\mathbf1,u^1,u^1)/\sqrt2
       -\rho_{\F}(\mathbf1,u^0,u^0)/\sqrt2$.
This fixes the auxiliary $i$ directly. It is independent of a numerical
comparison of blocks.

For implementation, remove the leftmost mode of a nontrivial PBW word,
choose a residue identity containing that mode with a nonzero coefficient,
apply all other modes by \eqref{eq:app-sca}, and solve for the requested
form. Memoize triples of words, ground labels, representation parameters,
and form labels. The ground tables terminate the reduction. The following
all-NS and NS--R--R descendant values are generated from the
displayed Ward identities. In particular, the global
superconformal seed uses the same BPZ forms; it needs no separate
odd three-point normalization.

\subsection{Sewing signs and spin-frame changes}
\label{app:sewing}

For states with the auxiliary factor first, the ordered enlarged
three-point form is
\begin{equation}
 \hat\rho(a_1\otimes x_1,a_2\otimes x_2,a_3\otimes x_3)
 =(-1)^{|a_1||x_1|+|x_2||a_3|}
   \rho_{\F}(a_1,a_2,a_3)\rho(x_1,x_2,x_3).
 \label{eq:app-enlarged-vertex}
\end{equation}
The first sign comes from the graded BPZ dual at infinity; the
second moves the middle physical state past the third auxiliary
state. An oscillator-frame change is performed on its arguments
before evaluating this form, not appended as another local phase.
Every edge uses the inverse of the full graded tensor Gram
\eqref{eq:app-product-metric}; no standalone
$(-1)^{\epsilon'_e}$ factor is attached.

For external states, extend the paper's definition of $K$ by placing
the ordered external states after the ordered internal edge-end pairs,
then count inversions into vertex order. External states are fixed,
not summed. A self-loop has two distinct slots but only one sewn edge.
The exact orders used in the release are listed in Sec.~\ref{app:graphs}.

All blocks in this appendix omit primary powers $q_e^{h_e}$.
The enlarged Ramond ground additionally has auxiliary weight $1/16$.
The omitted powers are restored only in a full correlator. These are
plumbing-plane powers, with no implicit $q^{-c/24}$ cylinder factor.

The mixed sign in the convolution is
\begin{equation}
 (-1)^{K(\epsilon+\epsilon')-K(\epsilon)-K(\epsilon')
        +\sum_e\epsilon_e\epsilon'_e
        +\sum_v(\epsilon_{v,1}\epsilon'_{v,1}
                +\epsilon_{v,2}\epsilon'_{v,3})}.
 \label{eq:app-convolution-sign}
\end{equation}
The first terms compare total and separate operator orderings,
the edge term comes from the tensor inverse Gram, and the last
terms come from \eqref{eq:app-enlarged-vertex}. Retain the $\eta_e$ as
formal variables with $\eta_e^2=1$ until recovery is complete.
The convolution is associative because its sign exponent is bilinear,
but its factors retain the displayed physical--auxiliary order.
Evaluating spin lifts before convolution loses the mixed parity signs.

Changing the local spin frame by $s_v=\pm1$ sends
$\eta_e\mapsto s_v\eta_es_w$ on an edge joining $v,w$.
For a closed block this changes its representative by the known scalar
$\prod_v s_v^{a_v}$ or $s_v^{f_v}$, as follows from the vertex parity
constraint. External states must be transformed too when present.
Two edge-sign assignments on a connected graph differ only by these
frame choices precisely when their products agree on every graph cycle.
One can check this by fixing frames along a spanning tree and comparing
the remaining edges. A self-loop's sign cannot be changed this way.
This criterion concerns spin lifts on a fixed NS/R-labelled channel;
it does not identify vertex-form signs or assign an Arf invariant.

\subsection{Primary scales and norm branches in the implementation}
\label{app:branches}

Normal ordering in the paper's embedding moves annihilators to the
right, with a minus sign for each fermion exchange. Its embedded
generators obey $(L_n^{(j)})^\dagger=L_{-n}^{(j)}$ under
\eqref{eq:app-adjoint}: the phases of $G^\dagger$ and
$\psi^\dagger$ cancel after graded reversal in $U_n$.
The code evaluates the displayed
$c^{(j)}$ and $h_n^{(j)}$ directly, without introducing square roots
of auxiliary Virasoro momenta.

Use the free-field realization in the paper's Appendix~A, with
$\{\eta_r,\psi_s\}=0$. In its oscillator basis,
$\eta_0$ interchanges the two grounds with coefficient $1/\sqrt2$,
giving $G_0=-iP\eta_0$ after the stated basis identification.
This oscillator realization constructs primaries and local actions;
the physical Gram matrix still uses the BPZ pairing above.

To specify the scale of these vectors, define the polynomial used
in their oscillator normalization:
\begin{equation}
 \ell(x,m)=2^{(1-(-1)^m)/16}
 \prod_{\substack{r,s\ge0,\ r+s<m\\r+s\equiv m\ (\mathrm{mod}\,2)}}
 (x+rb+s/b),\qquad m\ge0.
 \label{eq:app-ell}
\end{equation}
The empty product is one. For negative arguments in its integer slot,
\begin{equation}
 \ell(x,-m)=
 \begin{cases}
 \ell(Q-x,m),&m>0\text{ odd},\\
 (-1)^{m/2}\ell(Q-x,m),&m>0\text{ even}.
 \end{cases}
 \label{eq:app-ell-negative}
\end{equation}
For $n>0$, the implementation multiplies the NS oscillator product
in the paper's Appendix~A by
\begin{equation}
 (-1)^{n(2n-1)}2^{-2n}\ell(Q+2P,4n).
 \label{eq:app-ns-primary}
\end{equation}
The sign reverses the order of its $2n$ anticommuting $\chi$ modes;
the code places $\chi_{-1/2}$ leftmost. Below $v_n$ denotes this
rescaled vector, with $v_0$ unchanged and the same reflection rule.
In particular,
\begin{equation}
 v_{\pm1/2}=G_{-1/2}\phi_h+(Q/2\pm P)\psi_{-1/2}\phi_h,
 \qquad
 \|v_n(P)\|_{\rm BPZ}^2
 =\langle\!\langle v_n(P)|v_n(P)\rangle_{\eqref{eq:app-product-metric}}.
 \label{eq:app-ns-norm}
\end{equation}
The first expression suppresses the vacuum tensor factors. Expand
the vector in tensor PBW states, evaluate every term with
\eqref{eq:app-adjoint-pair} and \eqref{eq:app-product-metric}, and
cache the resulting bilinear norm. Use reflection for negative $n$.
Rescaling a primary would leave a fully
sewn block unchanged, but would change individual unnormalized
branching coefficients; the factor \eqref{eq:app-ns-primary} fixes
that scale.

The R primaries retain the oscillator construction and reflection
rule of the paper's Appendix~A. In that construction
$\chi_0(-P)=\psi_0+i\eta_0(P)$ on the ground space. Their current
BPZ norms are obtained by the same finite contraction:
\begin{equation}
 \|v_n^\alpha(P)\|_{\rm BPZ}^2
 =\langle\!\langle v_n^\alpha(P)|v_n^\alpha(P)
   \rangle_{\eqref{eq:app-product-metric}},\qquad \alpha=0,1.
 \label{eq:app-r-norms}
\end{equation}
Use $P\mapsto-P$, $n\mapsto-n$ for negative $n$,
including the reflected odd physical component. The closed
oscillator-frame norm formulas from the earlier convention are
not substituted for these BPZ norms.

In the paper's normalized branching coefficients, use one norm root
per module and branch at all its occurrences. Equivalently, use unnormalized vertices and
one inverse squared norm on each edge, as the code generally does.
This latter form avoids a norm-root branch choice altogether. An external
normalized state, if requested, still requires its explicitly chosen root.

\subsection{How the branching and splitting coefficients are fixed}
\label{app:branching-reproduction}

For the boundary values in the paper's Sec.~4, solve the finite
oscillator-to-PBW change of basis, then evaluate the enlarged local
form with the Ward identities above. Include both Ramond parities
and both $\eta$; the parity constraint determines $f$.
These are local three-point computations, not a physical block
supplied to the double-Virasoro calculation.

For the local $\mathbf V,\mathbb V$ coefficients defined there,
enumerate the pairs of Virasoro partition words at each specified
level, construct their oscillator vectors with the embedding, and
solve the finite change of basis for the physical $L_{\pm1}$ action.
Cache these local solutions. They are distinct from the recursive
evaluation of branching coefficients.
At $n_{\R}=1/4$, $L_{-1}$ also contains $v_{-3/4}$; at
$n_{\NS}=\pm1/2$ it mixes the two half-level branches. Reflection
gives the corresponding negative-label actions. These boundary couplings
must be retained.

Insert the cached actions into the two $L_1$ identities in the paper's
Sec.~4 and reduce their Virasoro descendants with
\eqref{eq:app-vir-ward}, normalized by
$\rho(\nu_1,\nu_2,\nu_3)=1$. Useful one-mode values are
\begin{align}
 \rho(L_{-k}\nu_1,\nu_2,\nu_3)&=h_1+kh_2-h_3,\nonumber\\
 \rho(\nu_1,L_{-k}\nu_2,\nu_3)&=(-1)^k(h_2+kh_3-h_1),\nonumber\\
 \rho(\nu_1,\nu_2,L_{-k}\nu_3)&=h_3+kh_2-h_1.
 \label{eq:app-single-vir-mode}
\end{align}
For repeated modes apply the Ward identity successively; the weight of
the descendant in the affected slot increases by its accumulated level.
The $L_0$ in the second identity is the \emph{physical} operator.
A double-Virasoro primary is generally not its eigenvector;
its $L_0$ action must also be expanded. Only $L_0^{(j)}$ acts on that
primary by $h_n^{(j)}$.

Evaluate unnormalized branching coefficients in increasing
\begin{equation}
 \lfloor|n_1|\rfloor+\lfloor|n_2|-\tfrac14\rfloor
                         +\lfloor|n_3|-\tfrac14\rfloor.
 \label{eq:app-branch-height}
\end{equation}
After inserting previously stored values, a bulk equation is one scalar
division. If its leading coefficient vanishes, use the second $L_1$
identity. The explicitly coupled boundary
labels require small simultaneous solves (at most four unanchored values
in the implementation). Store both $\eta$ right-hand sides, local action
vectors, weights, and one-leg Virasoro vertex products. Do not divide by
a zero Ward coefficient or discard the corresponding branch.
An unresolved pivot in both identities requires more precision or a
specified parameter limit, rather than an arbitrary replacement value.

All-NS branching on a general graph can be obtained by the same finite
oscillator/PBW construction and cached by its three labels and momenta.
For two identical theta vertices, an internal oscillator-frame
closed square is a useful shortcut, but it is not the definition
of the BPZ-normalized $\mathbf B_a$ here. Its phases must be
transported with the Gram and vertex basis change. The
convention-independent prescription is to construct each primary,
evaluate \eqref{eq:app-enlarged-vertex} with the displayed Ward
identities, and divide by the roots of its BPZ norms from
\eqref{eq:app-product-metric}. Cache the result under all three
branch labels and momenta. This also fixes relative phases between
different adjacent triples; independent square roots of a squared
coefficient would not.

\subsection{Insertion normalization and diagonal recovery}
\label{app:insertion}

Use the split channel and $\Theta$ action defined above. The
production code inserts $Q\Theta\psi$ in place of
$\Theta\psi$ at each new sphere, with the same scalar in its
auxiliary factor. The physical part of the inserted matrix element
is the unchanged physical Gram matrix;
contracting it between the two inverse physical Grams gives one inverse
Gram. This is why recovery yields the original physical block.

At $P=Q/2$ use the NS identity \emph{quotient}, with
$G_{-1/2}\mathbf1=0$. The primary expression in
\eqref{eq:app-ns-norm} then gives
$v_{1/2}=Q\psi_{-1/2}\mathbf1$. One must not obtain this insertion by
inverting the singular Gram matrix of an unreduced $h=0$ Verma module.
Use the two external Virasoro weights given in the paper's Sec.~5.
The allowed R branches obey
\begin{equation}
 n_R=n_L\pm\tfrac12,\quad\alpha_R=\alpha_L.
 \label{eq:app-insertion-branches}
\end{equation}
Evaluate the ground matrix elements in both orientations directly
from the oscillator primaries, the action
$\Theta u^{\mathsf a}=(-1)^{\mathsf a+1}u^{1-\mathsf a}$,
and the graded BPZ pairing \eqref{eq:app-product-metric}. Do not
identify the two orientations by an ungraded transpose. Apply
the splitting Ward identity above, or its
reversed version, using the same cached actions. A vanishing
middle pivot is currently treated as a nongeneric parameter requiring
a limit; the code does not silently regularize it.

Each split segment carries its own full graded tensor inverse Gram
before the insertion is contracted. In the fully expanded PBW sewing
\eqref{eq:split-enlarged-block}, multiply each inserted matrix element
by $Q$ for the production normalization. The punctured Virasoro blocks
are evaluated by the same $c$-recursion as the ordinary blocks.

Select equal \emph{combined} powers on each split pair. This is a
coefficient extraction, not a substitution $q_{e_L}=q_{e_R}$ and not
the requirement that the two Virasoro copies separately have equal
levels. The retained auxiliary operator acts as
\begin{equation}
 Q\Theta\psi_0\Psi_{-\mathsf A}u^{\mathsf a}
 =\frac{Q(-1)^{\mathsf a}}{\sqrt2}\Psi_{-\mathsf A}u^{\mathsf a}.
 \label{eq:app-diagonal-insertion}
\end{equation}
The nonzero-mode signs cancel because both $\Theta$ and $\psi_0$
are odd. In theta the auxiliary constant is $1-\eta_2\eta_3$ in
the ordinary calculation and
$Q(1+\eta_2\eta_3)/\sqrt2$ in the inserted one. On the required
cycle sector these act by $2$ and $\sqrt2Q$, respectively.
In general the corresponding scalar is
\begin{equation}
 \left(\prod_{C\text{ closed R cycle}}2\right)
 \left(\prod_{e\text{ split}}\frac Q{\sqrt2}\right).
 \label{eq:app-recovery-constant}
\end{equation}
Recover each physical coefficient by subtracting all positive-level
fermion convolutions with lower physical coefficients and dividing by
\eqref{eq:app-recovery-constant}. This is \eqref{eq:recovery} in the
production normalization. It is a restricted inverse on the required
cycle sector. No inverse of $1-\eta_2\eta_3$ on the full parity algebra
is used. Using canonical $\Theta\psi$ in both numerator and auxiliary
factor divides each insertion by $Q$ and leaves the recovered block
unchanged; the numerical $Q\Theta\psi$ convention requires $Q\ne0$.

\subsection{Virasoro recursion: pole branches and residue signs}
\label{app:ccy}

Every Virasoro three-point form has primary value one and uses
\eqref{eq:app-vir-ward}. Its internal primary has norm one. The block
omits all primary powers and has constant coefficient one. Recursion
varies the Virasoro central charge at fixed weights, independently in
each copy; it does not vary the physical $b$ or $P$ during that recursion.

For the pole on an edge of current weight $h$, take $r\ge2$, $s\ge1$,
with null level $rs$ in the allowed descendant domain. Put
\begin{align}
 x_\pm&=\frac{rs-1+2h\ \pm\sqrt{(r-s)^2+4(rs-1)h+4h^2}}{1-r^2},
 \nonumber\\
 x&=\begin{cases}x_+,&|x_+|\ge|x_-|,\\x_-,&|x_-|>|x_+|,
 \end{cases}
 \qquad c_{rs}(h)=13+6(x+x^{-1}).
 \label{eq:app-ccy-pole}
\end{align}
The square root is the principal complex root; a tie selects $x_+$.
Selecting the larger-magnitude root avoids cancellation in this
parameterization. This rule describes the implementation, rather than
an independently selected Liouville-momentum sign. If $x=0$ or poles
coincide, a generic-parameter limit is required.

The null vector is normalized by the coefficient of $L_{-1}^{rs}$ being
one. Its inverse-norm factor and the central-charge Jacobian give
\begin{align}
 A_{rs}&=\frac12
 \prod_{\substack{u=1-r,\ldots,r\ ;\ v=1-s,\ldots,s\\
                    (u,v)\ne(0,0),(r,s)}}
             (u\sqrt x+v/\sqrt x)^{-1},\nonumber\\
 -\frac{dc_{rs}}{dh}A_{rs}
 &=-\frac{12(x^2-1)x^{2rs-1}}
 {[(1-r^2)x^2-(1-s^2)]
  \displaystyle\prod_{\substack{u=1-r,\ldots,r\ ;\ v=1-s,\ldots,s\\
                    (u,v)\ne(0,0),(r,s)}}(ux+v)}.
 \label{eq:app-ccy-prefactor}
\end{align}
The denominator in the block is $c-c_{rs}(h)$, so the minus sign in
\eqref{eq:app-ccy-prefactor} must be retained. The code evaluates the
rational expression on the second line. It cancels explicit factors
$x-1$ and $x+1$ algebraically before numerical division; it does not
use a finite-difference derivative.

The ordered fusion polynomial can also be evaluated without spectator
momentum roots. For spectators of weights $h_a,h_b$, set, only in this
formula, $\lambda_a^2=x+2+x^{-1}-4h_a$ and
$\delta_{uv}=u^2x+2uv+v^2/x$. Over
$u=1-r,3-r,\ldots,r-1$ and $v=1-s,3-s,\ldots,s-1$, use
\begin{equation}
 P_{rs}(h_a,h_b)=
 (h_b-h_a)^{\mathbf1_{r,s\text{ both odd}}}
 \prod_{\substack{(u,v)\text{ in these ranges}\\u>0\ \text{or}\ (u=0,v>0)}}
 \frac{[4(h_b-h_a)+\delta_{uv}]^2-4\lambda_a^2\delta_{uv}}{16}.
 \label{eq:app-vir-fusion}
\end{equation}
This pairs opposite shifts before multiplication. For a null inserted in
slot $1,2,3$, the ordered spectator slots $(a,b)$ are $(3,2)$, $(3,1)$,
and $(1,2)$, respectively. Multiply the two incident fusion polynomials
by \eqref{eq:app-ccy-prefactor}; the child has this edge weight shifted
by $rs$ and central charge $c_{rs}$. Include the factor $q_e^{rs}$.
For a self-loop occupying slots $2,3$ of one sphere, with other weight
$h_1$ and loop weight $h$, the product is instead explicitly
\begin{equation}
 (-1)^{rs}P_{rs}(h_1,h+rs)P_{rs}(h_1,h).
 \label{eq:app-selfloop-residue}
\end{equation}
The shifted spectator and the displayed sign are essential for the
glasses channel.

The large-$c$ weight-dependent seed sews only $L_{-1}$ descendants.
Their Gram entry is $k!(2h)_k$, with $(x)_k=x(x+1)\cdots(x+k-1)$.
For complete reproducibility its local vertex is
\begin{align}
 &\rho(L_{-1}^i\nu_1,L_{-1}^j\nu_2,L_{-1}^k\nu_3)
 =(h_1-h_2-h_3+i-j+1-k)_j\nonumber\\
 &\quad\times\sum_{p=0}^{\min(i,k)}\binom ip
 \frac{k!}{(k-p)!}
 \left[\prod_{t=0}^{p-1}(2h_3+k-1-t)\right]
 (h_2+h_3-h_1)_{k-p}(h_1+h_2-h_3+p-k)_{i-p}.
 \label{eq:app-global-vertex}
\end{align}
Empty products equal one. Sew these vertices with $1/[k!(2h)_k]$
on each edge. The full seed also contains the Schottky vacuum factor
in Sec.~\ref{app:schottky}; retaining the global vertices alone is not
the higher-genus large-$c$ limit.

The code divides out the weight-independent Schottky factor while
propagating residues. Paths with the same accumulated edge shifts and
incoming central-charge pole are added before any global coefficient
is attached. Pole data, $A_{rs}$, fusion polynomials, inverse differences
of central charges, and local global vertices are cached at their full
parameter labels. The two Schottky factors are restored after the two
Virasoro copies and the branching sum have been combined. These are
reorganizations of the same recursion and do not alter its normalization.

\subsection{Schottky multipliers and the vacuum seed}
\label{app:schottky}

This subsection specifies the plumbing realization, so a Schottky
multiplier cannot be replaced by an independently chosen modulus.
Represent the three local-coordinate maps by
\begin{equation}
 C_1=\begin{pmatrix}0&1\\1&0\end{pmatrix},\qquad
 C_2=\begin{pmatrix}1&-1\\0&1\end{pmatrix},\qquad
 C_3=\begin{pmatrix}1&0\\0&1\end{pmatrix},
 \label{eq:app-coordinate-matrices}
\end{equation}
acting by $(az+b)/(cz+d)$. A directed edge from slot $i$ to slot $j$
has matrix
\begin{equation}
 M_{e,\to}=C_j^{-1}\begin{pmatrix}0&q_e\\1&0\end{pmatrix}C_i.
 \label{eq:app-edge-map}
\end{equation}
For a closed directed walk multiply the next matrix on the left.
Remove immediate backtracking, including between the last and first
edge. Primitive words are not positive powers of shorter words.
Keep one representative modulo cyclic rotation \emph{and inversion}.
This inverse-word identification fixes the exponent of the product below.

For a primitive matrix $M_\gamma$, let $T=\operatorname{tr}M_\gamma$
and $D=\det M_\gamma$. Its small multiplier is the formal series
\begin{equation}
 k_\gamma=\frac{1-\sqrt{1-4D/T^2}}{1+\sqrt{1-4D/T^2}},
 \qquad \sqrt{1-4D/T^2}=1+O(q),
 \qquad
 V_{\mathrm{Vir}}=\prod_{[\gamma]}\prod_{m=2}^{\infty}(1-k_\gamma^m)^{-1}.
 \label{eq:app-schottky}
\end{equation}
The constant of $T$ is nonzero for the primitive cusp words used here.
All inverses and roots in this construction are \emph{formal power
series at the plumbing cusp}. Numerically choosing the larger eigenvalue
after evaluating a truncated matrix is not the prescription.
The equivalent series is
$k_\gamma=\sum_{l\ge1}\frac1{l+1}\binom{2l}{l}(D/T^2)^l$.
It is implemented in exact rational arithmetic before evaluation at
representation parameters. There is no $m=1$ oscillator in
$V_{\mathrm{Vir}}$; it belongs to the global sector already sewn.

If a word visits edge $e$ a total of $v_e$ times, its first bosonic
oscillator starts with $\prod_e q_e^{2v_e}$. Therefore expand its
unit trace factor only in the domain remaining after this monomial.
For total cutoffs subtract $2\sum_e v_e$; for independent edge cutoffs
subtract $2v_e$ separately. This is an exact truncation rule. It is
particularly useful when the target is a box of edge levels.

The punctured theta numerator with a physical identity at the inserted
sphere uses the same surface vacuum factor after $q_e=q_{e_L}q_{e_R}$.
On a general expanded graph the same matrix prescription may instead
be applied directly. In either approach the starting monomials and
the final diagonal extraction are retained before truncation.

\subsection{Conventions for the independent all-NS reference}
\label{app:ns-reference}

The comparison with SCA $c$-recursion uses fixed physical NS weights,
not the two Virasoro weights. Nulls have level $rs/2$, with
$r\ge2$, $s\ge1$, $r+s$ even. Replace the pole formula by
\begin{align}
 x_\pm&=\frac{4h+rs-1\ \pm
 \sqrt{16h^2+8(rs-1)h+(r-s)^2}}{1-r^2},\nonumber\\
 c_{rs}^{\NS}(h)&=\frac{15}{2}+3(x+x^{-1}),
 \qquad |x|=\max(|x_+|,|x_-|),
 \label{eq:app-ns-pole}
\end{align}
with the same principal root and tie rule. The prefactor is
$-(dc_{rs}^{\NS}/dh)A_{rs}^{\NS}$, where
\begin{equation}
 A_{rs}^{\NS}=\frac12
 \prod_{\substack{u=1-r,\ldots,r\ ;\ v=1-s,\ldots,s\\
  u+v\ \mathrm{even},\ (u,v)\ne(0,0),(r,s)}}
 \frac{\sqrt2}{u\sqrt x+v/\sqrt x}.
 \label{eq:app-ns-prefactor}
\end{equation}
Normalize an even null by the coefficient of $L_{-1}^{rs/2}$ being
one, and an odd null by that of $L_{-1}^{(rs-1)/2}G_{-1/2}$ being one.
Use the ordered spectators of Sec.~\ref{app:ccy}. In the fusion formula
replace $4$ by $8$ in $\lambda_a^2$ and the weight difference, and
$16$ by $64$ in the denominator. Keep only shifts with
\begin{equation}
 u+v-r-s\equiv
 \begin{cases}2\pmod4,&a_v=0,\\0\pmod4,&a_v=1.
 \end{cases}
 \label{eq:app-ns-fusion-selection}
\end{equation}
Opposite shifts are paired as before; the zero shift, when present in
this subset, contributes $h_b-h_a$.

An odd null changes the intrinsic parity of the shifted primary.
Express the child again in the even-primary convention using the
graded BPZ pairing and the all-NS Ward identities
\eqref{eq:app-ns-g-ward}. Compute its slot incidence factors from
those identities at both ends, together with
$(-1)^{K(\epsilon)-K(\epsilon_{\mathrm{child}})}$;
the old frame's fixed three-slot phase table does not apply.
Here $\epsilon_e=2\,\mathrm{level}_e\bmod2$ before the null shift,
and the child's edge parity changes if $rs$ is odd. The fusion label
at each incident vertex is the sum of its three current edge parities.
This specifies the parity transport, rather than treating every
residue as a scalar bosonic shift.

The global NS basis is $L_{-1}^kG_{-1/2}^a\phi_h$, $a=0,1$,
with BPZ Gram entry $i^a k!(2h)_{k+a}$. Its vertex is
\eqref{eq:app-global-vertex} with $h_i$ replaced by $h_i+a_i/2$,
multiplied by the corresponding value generated by
\eqref{eq:app-ns-g-ward}.
Include the physical $(-1)^K$ when sewing. The vacuum factor is
\begin{equation}
 V_{\NS}=V_{\mathrm{Vir}}
  \prod_{[\gamma]}^{\star}\prod_{m=2}^{\infty}
             (1-u_\gamma k_\gamma^{m-1}),
 \qquad u_\gamma\star u_\gamma=k_\gamma.
 \label{eq:app-ns-vacuum}
\end{equation}
The first fermionic oscillator is at $3/2$, not $1/2$; the latter is
already in the global basis. The sign of the half-multiplier $u_\gamma$
is fixed by the spin lift, not by a principal numerical square root.
Sew the global factor on the left of $V_{\NS}$ using
\eqref{eq:app-convolution-sign}; its local middle-to-ket term generally survives
on the tetrahedron.

For clarity the following gives the exact half-multiplier prescription
used for the two reference graphs. Let $d_e=\det M_{e,\to}/q_e$,
which is $\pm1$. On theta $d_e=(-1,-1,-1)$, and on the ordered
tetrahedron of \eqref{eq:mercedes-vertices}
$d_e=(-1,1,1,-1,-1,-1)$. Assign edge phases
$(0,0,0)$ and $(0,0,0,0,2,0)$, respectively, in units of $\pi/2$.
For each traversal add its edge phase, and add $1$ for a forward or
$3$ for a reverse traversal if $d_e=-1$. Denote this sum by $p$ just
within this prescription.

There is a necessary conversion from scalar spin lifts to the formal
parity algebra. For a cycle parity $\epsilon$, define $\tau(\epsilon)$
modulo four by $\tau(0)=0$ and
\begin{equation}
 i^{\tau(\epsilon)+\tau(\epsilon')-\tau(\epsilon+\epsilon')}
 =(-1)^{K(\epsilon+\epsilon')-K(\epsilon)-K(\epsilon')
            +\sum_e\epsilon_e\epsilon'_e
            +\sum_v(\epsilon_{v,1}\epsilon'_{v,1}
                         +\epsilon_{v,2}\epsilon'_{v,3})}.
 \label{eq:app-cycle-phase}
\end{equation}
On theta take $\tau=0$ on each of the basis cycles $\{1,2\}$ and
$\{1,3\}$. On the tetrahedron take $\tau=0,2,0$ on the cycles
$\{1,2,4\}$, $\{2,3,5\}$, and $\{1,3,6\}$, respectively.
Equation~\eqref{eq:app-cycle-phase} extends these
choices to every sum of basis cycles. These are Schottky
half-multiplier branch choices in the current spin frame, not a
conversion factor applied to a completed block. If a primitive word has visits
$v_e$ and parity $\epsilon_e=v_e\bmod2$, then
\begin{equation}
 u_\gamma=
 i^{p-\tau(\epsilon)}\,
 \frac{\prod_e q_e^{v_e/2}\eta_e^{\epsilon_e}}{T}
 \sum_{l\ge0}\frac1{l+1}\binom{2l}{l}(D/T^2)^l.
 \label{eq:app-half-multiplier}
\end{equation}
The primitive words, $T$, and $D$ use \eqref{eq:app-edge-map} with
the ordered edges in Sec.~\ref{app:graphs}. This formula specifies all
phases for the saved reference computations. Its convolution square
is the multiplier in \eqref{eq:app-schottky}. The product in
\eqref{eq:app-ns-vacuum} starts at half-level monomial $3v_e$;
integer-power corrections therefore need only the remaining cutoff
after $\lceil3v_e/2\rceil$ on each edge, or
$\lceil3\sum_e v_e/2\rceil$ at total cutoff.

The batched SCA reference groups residue paths with the same accumulated
null shifts and incoming pole separately for each requested output parity.
At fixed output parity their residues do not depend on the remaining
even levels. The superglobal coefficients are attached after this sum,
and the NS vacuum factor is restored once. The original coefficient-wise
recursion remains available for the directed comparison; neither uses
the double-Virasoro numerator as input.

\subsection{Ordered graphs, cutoffs, and numerical precision}
\label{app:graphs}

The following ordered slots specify the graphs used in the release.
An edge endpoint is written $(v,i)$ for vertex $v$, slot $i$.
Number vertices in the order of the rows shown, and edges in numerical
order. The slots and oriented endpoint pairs are
\begin{equation}
\begin{array}{c|c|l}
\text{channel}&\text{vertex slots }(\infty,1,0)&\text{edge endpoints in order}\\\hline
\theta&(1,2,3),(1,2,3)&
 ((1,1),(2,1)),\ ((1,2),(2,2)),\ ((1,3),(2,3))\\
\mathrm{glasses}&(1,2,2),(1,3,3)&
 ((1,1),(2,1)),\ ((1,2),(1,3)),\ ((2,2),(2,3))
\end{array}
 \label{eq:app-genus-two-order}
\end{equation}
For the tetrahedron the slots are those of
\eqref{eq:mercedes-vertices}; its six pairs are
\begin{equation}
 \bigl((1,1),(2,1)\bigr),\quad
 \bigl((1,2),(3,1)\bigr),\quad
 \bigl((1,3),(4,1)\bigr),\quad
 \bigl((2,2),(3,3)\bigr),\quad
 \bigl((3,2),(4,3)\bigr),\quad
 \bigl((4,2),(2,3)\bigr).
 \label{eq:app-tetrahedron-order}
\end{equation}
In the NS--R--R theta test edge 1 is NS. In the glasses tests edge 1
is NS and each handle may be NS or R. In the Ramond tetrahedron test
edges $1,2,3$ are NS and $4,5,6$ are R. The other tetrahedron reference
test is all-NS. These assignments do not include its inequivalent
four-edge Ramond cycle.

A total cutoff $L$ means $\sum_e\mathrm{level}_e\le L$; an
independent-edge cutoff means $\mathrm{level}_e\le L$ on every edge.
These levels exclude primary weights, and are half-integral on NS
edges and integral on R edges. In a branch, subtract $2n^2$ or
$2n^2-1/8$ before allocating descendant levels between the two
Virasoro copies. On an inserted edge a target retained at original
level $l$ has \emph{both} segment powers equal to $l$. Equivalently
an intermediate split monomial is needed in the downward closure when
its contribution, charged as the maximum of the two segment levels,
fits the original physical cutoff. Intermediate unequal segment
levels are necessary. A total bound obtained by adding the two segment
levels would impose the wrong physical domain.

The general-graph output field \texttt{level2[e]} equals twice the
physical exponent of $q_e$ in every sector. The specialized theta
engine internally stores $(2l_1,l_{2,L},l_{2,R},2l_3)$; its diagonal
has $l_{2,L}=l_{2,R}$. Internal branch labels are $4n$ in both engines.
A parity mask is $\sum_e2^{e-1}\epsilon_e$, with edge 1 the lowest
bit, and two masks are added by bitwise XOR. Output complex components
are stored as decimal real and imaginary strings. Missing components
are zero; they are included when comparing coefficient tables.
The present graph storage has eight internal edge slots, including
inserted segments. Local oscillator storage allows bosonic mode indices
below 64 and 64 fermion occupations, with boson occupation numbers at most
255. These implementation bounds are checked and are not genus or
level bounds on the mathematical construction. The supplied frontend
enumerates the named graphs above rather than accepting an arbitrary
plumbing graph from a file.

The default generic point is
\begin{equation}
 b=\frac75,\qquad
 (P_1,P_2,P_3)=\left(\frac{11}{23},\frac{13}{29},\frac{17}{31}\right),
 \qquad
 (P_4,P_5,P_6)=\left(\frac{19}{37},\frac{23}{41},\frac{29}{43}\right)
 \label{eq:app-test-parameters}
\end{equation}
when six edges are present. Supply these as rational strings, not
rounded machine numbers. Complex inputs use \texttt{real,imaginary},
with a rational or decimal string for each part. The current theta CLI
requires real $b\ne0,\pm1$; complex momenta are supported.
The graph frontend uses even intrinsic NS primary parities. The
specialized theta option \texttt{--p 1} keeps the same local numerical
forms but reverses the outgoing NS parity bit. In that interface
\texttt{--f} is the relative descendant/ground-parity label; the
absolute vertex parity is $f+p\bmod2$. The formulas in this appendix
and the graph tests take $p=0$.

Machine precision uses \texttt{std::complex<double>}. At requested
precision $D\ge30$ decimal digits, MPC uses
$\operatorname{round}((D+1)\log_2 10)$ bits and round-to-nearest
for both components. Square roots of numerical complex parameters use
the principal branch, positive on the positive real axis. The formal
Schottky roots, transported contour roots, and reflected momentum
choices above override any attempt to take unrelated pointwise roots.
Integer powers are evaluated by multiplication, including reciprocals
for negative powers. At a point where primary powers are restored,
$q^h=\exp(h\operatorname{Log}q)$ uses the principal logarithm in the
pointwise adapter. Analytic continuation around a loop requires
transporting that logarithm and the spin frame; the pointwise adapter
does not perform such transport automatically.

MPC arithmetic does not make every auxiliary decision multiprecision.
The local oscillator span solver scales columns, uses machine LAPACK
to select independent rows and build an LU preconditioner, then refines
residuals with the original multiprecision columns. It checks all rows,
including unselected ones. Its Hermitian Euclidean norms are solely
conditioning diagnostics and do not change the bilinear BPZ form.
The rank threshold is $10^{-11}$; refinement stops at relative error
$10^{-\max(20,D-15)}$, with at most 30 iterations. The full-row
acceptance threshold is $10^{-\max(15,D-20)}$ (machine: $10^{-8}$).
The small recursive branching systems are eliminated directly in the
chosen scalar precision.

Some sparse oscillator operations discard coefficients below
$10^{-\max(20,D-20)}$ (machine: $10^{-13}$). A branching pivot is
treated as unresolved when its magnitude relative to the sum of the
magnitudes of its contributions is at most
$10^{-\max(15,D-20)}$ (machine: $10^{-11}$). The Virasoro pole
separation threshold, relative to the larger central charge or one,
is $10^{10-D}$ (machine: $64\,2^{-52}$). The independent NS reference
also has an absolute pole-separation guard $10^{-25}$.
These are implementation thresholds, not analytic criteria for
vanishing. High precision can reduce cancellation error but does not
define a degenerate-module or confluent-pole limit. Near a rank failure
or an unresolved pole a separate limit or a different numerical
linear algebra routine is needed.

For comparison the reported error of a coefficient is
\begin{equation}
 \frac{|\mathbf F_{\mathrm{DV}}-\mathbf F_{\mathrm{reference}}|}
 {\max(1,|\mathbf F_{\mathrm{DV}}|,|\mathbf F_{\mathrm{reference}}|)}.
 \label{eq:app-error}
\end{equation}
The release threshold is $10^{-18}$, with maximum absolute errors
reported separately. Numerical working precision is not a statement
that every final coefficient has that many accurate digits. The
specialized recovery code additionally checks its cycle-sector residual
against $10^{-8}$ unless \texttt{--sector-policy record} was explicitly
selected; the latter records rather than suppresses this discrepancy.
Its ordinary-sector projection and the unprojected residual are recorded
as implementation choices. A failed comparison is retained, and the
graph runner repeats only that case with 20 additional digits in both
methods. It does not change the original failure into a passing result.

\subsection{Nonchiral data and identity normalization}
\label{app:nonchiral-conventions}

A partition function also requires a spectrum, a spectral measure,
and three-point coefficients. These are independent of the chiral
coefficient algorithm. Appendix~\ref{app:partition-test} specifies
these inputs for super-Liouville theory, using $C_a$ and $C_{f,\eta}$,
and writes the partition functions directly in the paper's blocks.
The BPZ factors are included in the returned block, and the
partition driver evaluates it at the chosen tube signs. The identity limit
in \eqref{eq:partition-identity-residue} fixes the relative NS/R
normalization. The bilinear BPZ definition is retained; complex
conjugation is imposed only on the physical diagonal spectrum.

\subsection{Reproduction commands and the scope of the checks}
\label{app:reproduce}

The release \path{C++/release/section6-superconformal-blocks.tar.gz}
includes file hashes in \path{SOURCE_MANIFEST.json}. It requires a
C++17 compiler, GMP, MPFR, MPC, and LAPACK/BLAS (Accelerate on macOS),
with Python 3.10 or later for the launcher and decimal comparisons.
Build from the extracted root:
\begin{verbatim}
make -C C++ section6
make -C C++ check-ns-forward check-graph-schottky
\end{verbatim}
The second command checks the batched NS recursion against its
coefficient-wise form and checks the Schottky truncation with exact
rational coefficients. Local physical, auxiliary, and embedded-Virasoro
Ward checks are also retained in the full repository under
\path{C++/experiments/coherent_conventions_2026-09-22/}; their
\path{README.md}, source, and JSON results specify their finite scope.
Those checks include the numerical ground-basis conversion described
in Sec.~\ref{app:pbw}.
The local suite was rebuilt and repeated while preparing this appendix
on September 23, 2026: all 8,238 checked values passed at 40 digits,
with maximum scaled discrepancy $3.01\times10^{-31}$ in the embedded
Virasoro Ward check. The fresh results and source hashes are retained in
\path{C++/results/conventions_appendix_2026-09-23/}.

\noindent\begin{minipage}{\linewidth}
For a fresh theta calculation with the default parameters, use one of
the following and select a previously unused output filename:
\begin{verbatim}
C++/bin/ramond --mode ordinary --level 5 \
  --truncation per-edge --dps 40 --json /tmp/theta-ordinary.json
C++/bin/ramond --mode inserted --level 5 \
  --truncation per-edge --dps 40 --json /tmp/theta-inserted.json
\end{verbatim}
\end{minipage}

The ordinary mode takes $\eta'=\eta$ and the inserted mode takes
$\eta'=-\eta$. The options \texttt{--f}, \texttt{--eta},
\texttt{--p}, \texttt{--b}, and \texttt{--P1} through
\texttt{--P3} select the remaining data. The separate executable
\texttt{C++/bin/pbw} supplies the physical reference. Machine precision
is \texttt{--dps 0} for these theta executables; the general-graph
frontend currently requires MPC precision of at least 30 digits.

To reproduce the graph PBW comparisons and the all-NS recursion
comparisons, use fresh output directories:
\begin{verbatim}
python3 C++/tools/section6.py run --suite theta-pbw \
  --output /tmp/section6-theta-pbw
python3 C++/tools/section6.py run --suite graph-pbw \
  --output /tmp/section6-graph-pbw
python3 C++/tools/section6.py run --suite theta-ns \
  --output /tmp/section6-theta-ns
python3 C++/tools/section6.py run --suite tetrahedron-ns \
  --output /tmp/section6-tetrahedron-ns
\end{verbatim}
These commands perform full computations. The genus-two PBW comparisons
use independent edge level 5; the triangular-R tetrahedron PBW comparison
uses total level 5. The all-NS theta reference uses independent edge
level 10, and the all-NS tetrahedron reference uses total level 10.
Every allowed local form/sign choice in those tests is retained, along
with all parity components and spin-lift evaluations. The graph PBW
suite also checks uninserted odd-loop cancellation, recovery, forward
convolution, and unequal split powers. In particular, it does not test
only a physical diagonal while claiming a full split-series comparison.

The coefficient engine itself is selected by
\texttt{C++/bin/graph\_blocks --method dv}; it evaluates physical
coefficients by double Virasoro without computing a physical-PBW
reference. Use \texttt{--method pbw-check} for the comparison and
\texttt{--method ns} for the independent all-NS reference. The graph
and all its form/sign cases are recorded in \path{metadata.json};
\path{coefficients.jsonl} contains the coefficient rows and
\path{summary.json} is written after completion. These files, the
source hashes, commands, precision, and timings are retained by the
launcher. A successful process exit without a passing comparison is
not reported as agreement.

The consolidated results are in
\path{C++/results/section6_2026-09-23/validation_summary.json}.
One glasses case failed the $10^{-18}$ criterion at 40 digits and
passed after both methods were repeated at 60 digits; the original
failure is preserved. The all-NS tetrahedron level-10 evidence is a
complete previously saved computation, re-compared and hash-checked,
with fresh production prefixes through total level 3 (double Virasoro)
and total level 6 (batched NS recursion). It is not described as a new
level-10 timing run. The independent level-10 all-NS theta calculation
is fresh. These numerical checks are at generic momenta, and do not
prove an arbitrary-genus statement by numerical enumeration.

The current nonchiral conventions are implemented in
\path{C++/include/scblocks/partition_sewing.hpp}. Build
\texttt{make -C C++ partition} and use \texttt{C++/bin/partition};
\path{C++/PARTITION.md} gives the inputs and complete commands.
The three-point formulas and the block expressions are given in
Appendix~\ref{app:partition-test}. The corrected fresh results,
including the geometric spin checks, are retained in
\path{C++/results/partition_paper_pairing_2026-09-24/}.
The independent identity-residue tests remain in
\path{C++/results/partition_normalization_resolution_2026-09-24/}.
Historical Python assemblers are retained to reproduce their
original data, not as the current coefficient convention.
